% 专题三：函数与方程
\section{专题三：函数与方程}

\begin{learnobjective}
\begin{itemize}
  \item 说出方程与函数的关系，会解简单的一元二次方程；
  \item 会用零点存在定理判断方程根的存在性；
  \item 会用极值判断方程解的个数；
  \item 会用数形结合处理简单的超越方程。
\end{itemize}
\end{learnobjective}

\begin{preview}
\subsection{导语：同一个问题，五种看法}
方程与函数是同一件事的两种说法：方程 $f(x)=0$ 的解，就是函数 $y=f(x)$ 的零点。
同一个方程问题，用不同的工具会有不同的理解：
直接求解、画图看交点、用零点存在定理判断有没有根、
用极值判断有几个根、用数形结合处理超越方程。
本专题把三个“种子问题”用五种视角各做一遍。

\subsection{种子问题}
先自己做，再看后面的五种视角。
\begin{enumerate}
  \item 解方程 $x^2-2x-3=0$。
  \item 判断方程 $x^3-x-1=0$ 有没有实数根，有几个。
  \item 判断方程 $2^x=x+2$ 的实数解个数。
\end{enumerate}
\end{preview}

\begin{explore}
\subsection{视角一：来源——方程从哪来}

\key{要点}
\begin{itemize}
  \item 方程是“未知量满足的等量关系”；解方程就是求使等式成立的未知量。
  \item 令 $y=f(x)$，则方程 $f(x)=0$ 的解就是函数 $y=f(x)$ 的\key{零点}。
  \item 三个二次：二次函数、一元二次方程、一元二次不等式是同一件事的三种说法。
\end{itemize}

\begin{example}
解方程 $x^2-2x-3=0$。
\end{example}
\begin{solution}
因式分解：$(x-3)(x+1)=0$，故 $x=3$ 或 $x=-1$。
对应函数 $y=x^2-2x-3$ 的零点为 $-1$ 和 $3$。
\key{（回头看：种子问题 1）}
\end{solution}

\begin{example}
解方程 $2x^2-4x-3=0$。
\end{example}
\begin{solution}
求根公式：$x=\dfrac{4\pm\sqrt{16+24}}{4}=\dfrac{4\pm2\sqrt{10}}{4}=\dfrac{2\pm\sqrt{10}}{2}$。
\end{solution}

\begin{problem}
解方程 $x^2-5x+6=0$，并写出对应函数的零点。
\end{problem}
\begin{solution}
$(x-2)(x-3)=0$，故 $x=2$ 或 $x=3$；对应函数的零点为 $2$ 和 $3$。
\end{solution}

\subsection{视角二：证明——根的存在性}

\key{要点}
\begin{itemize}
  \item \key{零点存在定理}：若函数 $y=f(x)$ 在闭区间 $[a,b]$ 上的图象是连续不断的，且 $f(a)f(b)<0$，则 $f$ 在 $(a,b)$ 内至少有一个零点。
  \item 它只保证“有”，不保证“只有一个”；要判断个数，需结合单调性或极值。
  \item \key{二分法}：反复把含零点的区间对半分，逐步逼近零点，用于求近似解。
\end{itemize}

\begin{example}
证明：方程 $x^3-x-1=0$ 在 $(1,2)$ 内有实数根。
\end{example}
\begin{solution}
设 $f(x)=x^3-x-1$，则 $f(1)=-1<0$，$f(2)=5>0$，且 $f$ 是多项式函数（连续），
由零点存在定理，$f$ 在 $(1,2)$ 内至少有一个零点，即方程在 $(1,2)$ 内有实数根。
\key{（回头看：种子问题 2 的存在性）}
\end{solution}

\begin{problem}
用二分法求方程 $x^3-x-1=0$ 的近似解，要求区间长度小于 $0.1$。
\end{problem}
\begin{solution}
零点在 $(1,2)$。$f(1.5)=0.875>0$，零点在 $(1,1.5)$；
$f(1.25)=-0.297<0$，零点在 $(1.25,1.5)$；
$f(1.375)=0.225>0$，零点在 $(1.25,1.375)$；
$f(1.3125)=-0.052<0$，零点在 $(1.3125,1.375)$，区间长度 $0.0625<0.1$。
近似解可取 $x\approx1.34$。
\end{solution}

\subsection{视角三：函数——方程的解就是函数的零点}

\key{要点}
\begin{itemize}
  \item 方程 $f(x)=0$ 的解 $=$ 函数 $y=f(x)$ 的零点 $=$ 图象与 $x$ 轴交点的横坐标。
  \item 方程 $f(x)=g(x)$ 的解 $=$ 两个函数图象交点的横坐标。
\end{itemize}

\begin{example}
用图象判断方程 $x^3-x-1=0$ 的实数解个数。
\end{example}
\begin{solution}
设 $f(x)=x^3-x-1$，$f'(x)=3x^2-1$，令 $f'(x)=0$ 得 $x=\pm\dfrac{1}{\sqrt3}$。
$f$ 在 $x=-\dfrac{1}{\sqrt3}$ 处取极大值 $\dfrac{2\sqrt3}{9}-1<0$；
在 $x=\dfrac{1}{\sqrt3}$ 处取极小值 $-\dfrac{2\sqrt3}{9}-1<0$。
极大值小于 $0$，且 $x\to+\infty$ 时 $f(x)\to+\infty$，
故图象与 $x$ 轴只有一个交点，方程只有一个实数解。
\key{（回头看：种子问题 2）}
\end{solution}

\begin{problem}
求函数 $f(x)=x^3-3x$ 的极值，并据此判断方程 $x^3-3x=1$ 的实数解个数。
\end{problem}
\begin{solution}
$f'(x)=3x^2-3=3(x-1)(x+1)$。$x=-1$ 处取极大值 $f(-1)=2$，$x=1$ 处取极小值 $f(1)=-2$。
方程 $x^3-3x=1$ 即 $f(x)=1$。因 $-2<1<2$，直线 $y=1$ 与 $y=f(x)$ 有三个交点，
故方程有三个实数解。
\end{solution}

\subsection{视角四：极值——用极值判断解的个数}

\key{要点}
\begin{itemize}
  \item 方程 $f(x)=k$ 的解的个数 $=$ 直线 $y=k$ 与曲线 $y=f(x)$ 的交点个数。
  \item 先求 $f$ 的单调区间与极值，再比较 $k$ 与极值的大小，即可判断交点个数。
\end{itemize}

\begin{example}
讨论方程 $x^3-3x+k=0$ 的实数解个数。
\end{example}
\begin{solution}
设 $f(x)=x^3-3x+k$，$f'(x)=3(x-1)(x+1)$。极大值 $f(-1)=2+k$，极小值 $f(1)=-2+k$。
\begin{itemize}
  \item $k<-2$ 或 $k>2$：极大值与极小值同号，方程有一个实数解；
  \item $k=\pm2$：极值点处相切，方程有两个实数解；
  \item $-2<k<2$：方程有三个实数解。
\end{itemize}
\key{（回头看：种子问题 2 的推广）}
\end{solution}

\begin{problem}
讨论方程 $x^3-6x+1=0$ 的实数解个数。
\end{problem}
\begin{solution}
$f(x)=x^3-6x+1$，$f'(x)=3x^2-6=3(x-\sqrt2)(x+\sqrt2)$。
极大值 $f(-\sqrt2)=4\sqrt2+1>0$，极小值 $f(\sqrt2)=-4\sqrt2+1<0$，
故方程有三个实数解。
\end{solution}

\subsection{视角五：几何——数形结合}

\key{要点}
\begin{itemize}
  \item 把方程 $g(x)=h(x)$ 的解看成两个函数 $y=g(x)$、$y=h(x)$ 图象交点的横坐标。
  \item 对含有指数、对数的超越方程，通常无法直接求解，数形结合是判断解的个数的主要工具。
\end{itemize}

\begin{example}
判断方程 $2^x=x+2$ 的实数解个数。
\end{example}
\begin{solution}
设 $f(x)=2^x-x-2$，$f'(x)=2^x\ln2-1$，令 $f'(x)=0$ 得 $x=-\log_2(\ln2)\approx0.53$。
$f$ 在 $(-\infty,-\log_2(\ln2))$ 上递减，在 $(-\log_2(\ln2),+\infty)$ 上递增，
极小值 $f(-\log_2(\ln2))\approx-1.09<0$。
又 $\lim\limits_{x\to-\infty}f(x)=+\infty$，$f(2)=0$，$\lim\limits_{x\to+\infty}f(x)=+\infty$，
故图象与 $x$ 轴恰有两个交点：一个为 $x=2$，另一个在 $(-2,0)$ 内
（因 $f(-2)=0.25>0$，$f(0)=-1<0$）。
即方程有两个实数解。
\key{（回头看：种子问题 3）}
\end{solution}

\begin{problem}
判断方程 $\ln x=x-2$ 的实数解个数。
\end{problem}
\begin{solution}
设 $f(x)=\ln x-x+2$（$x>0$），$f'(x)=\dfrac1x-1$，令 $f'(x)=0$ 得 $x=1$。
$f$ 在 $(0,1)$ 上递增，在 $(1,+\infty)$ 上递减，极大值 $f(1)=1>0$。
又 $x\to0^+$ 时 $f(x)\to-\infty$，$x\to+\infty$ 时 $f(x)\to-\infty$，
故图象与 $x$ 轴有两个交点，方程有两个实数解。
\end{solution}
\end{explore}

\begin{inclasstest}
\subsection{变式训练}
只改一个条件，重新做一遍。
\begin{enumerate}
  \item 判断方程 $x^3-3x+1=0$ 在 $(0,1)$ 内有没有实数根，并说明理由。
  \item 用二分法求方程 $x^3-3x+1=0$ 在 $(0,1)$ 内的近似解（区间长度小于 $0.2$）。
  \item 讨论方程 $x^2-2x+k=0$ 的实数解个数。
  \item 判断方程 $\ln x=x-2$ 的实数解个数。
\end{enumerate}
\end{inclasstest}

\begin{afterclass}
\subsection{知识来源对照表}
\begin{center}
\footnotesize
\setlength{\tabcolsep}{4pt}
\renewcommand{\arraystretch}{1.35}
\begin{tabular}{p{3.4cm}p{4.0cm}p{5.6cm}}
\toprule
结论 / 方法 & 最早从哪来 & 后面从哪个角度看 \\
\midrule
三个二次的关系 & 视角一（来源） & 视角三（函数）、视角五（几何） \\
零点存在定理 & 视角二（证明） & 视角三（函数）、视角五（几何） \\
零点即图象交点 & 视角三（函数） & 视角五（几何） \\
用极值判断解的个数 & 视角四（极值） & 视角三（函数） \\
二分法 & 视角二（证明） & 视角五（几何） \\
数形结合 & 视角五（几何） & 视角三（函数） \\
\bottomrule
\end{tabular}
\end{center}

\key{下一步}：把“函数与方程”迁移到别的主题（如极值），
同样用“种子问题 + 五个视角”重做一遍。
\end{afterclass}
